\chapter{Business rules}\label{ch:rules}
\chapterstatus{complete}{2026-11-17}{\src{src/domain}, \src{src/server/services}, \src{prisma/schema.prisma}}

The rules Slotty enforces today and where each one lives. Pure rules go in \code{src/domain/*} and return a \code{RuleResult} whose message is shown to the user. Services enforce them with \code{assertRule}. When you change a rule, update this chapter, the domain function and its unit test together.

\section{Demo policy}

One \code{DemoPolicy} per assignment (\src{prisma/schema.prisma}). Staff set it on the assignment form.

{\tablefont
\begin{xltabular}{\linewidth}{@{}P{30mm} P{16mm} L@{}}
  \toprule
  \thead{Field} & \thead{Default} & \thead{Meaning}\\
  \midrule
  \endhead
  \code{windowStart} / \code{windowEnd} & --- & Demos can only happen inside this window. Slots are clipped to it.\\
  \code{slotDurationMin} & --- & Length of each slot (5--480).\\
  \code{bufferMin} & 0 & Break between generated slots.\\
  \code{capacityPerSlot} & 1 & Students per slot (group demos if > 1). Copied onto each slot when it's created.\\
  \code{bookingOpensAt} & none & Students can't book before this time.\\
  \code{freezeHours} & 12 & No booking, cancelling or rescheduling within this many hours of the slot.\\
  \code{maxReschedules} & 2 & \textbf{Changes} allowed per student per assignment. Reschedules and self-cancellations both count (see change budget (\cref{sec:change-budget})).\\
  \code{allowStudentCancel} & true & Whether students can cancel on their own.\\
  \code{openAnnouncedAt} & --- & Internal: when students were told booking is open. The worker announces at \code{bookingOpensAt} if it's in the future.\\
  \bottomrule
\end{xltabular}}

\section[Booking rules]{Booking rules (\src{src/domain/booking-rules.ts})}

{\tablefont
\begin{xltabular}{\linewidth}{@{}P{30mm} P{30mm} L@{}}
  \toprule
  \thead{Rule} & \thead{Enforced in} & \thead{User message}\\
  \midrule
  \endhead
  Assignment must be \code{PUBLISHED} to book, reschedule or cancel (or \code{CLOSED} with a late-booking allowance) & \code{canBookSlot}, \code{canCancelBooking} & ``This assignment is not open for booking.'' / ``Booking is closed, so you can't cancel here. Ask your TA if you can't attend.''\\
  Slot must be \code{PUBLISHED} & \code{canBookSlot} & ``This slot is not available.''\\
  Not before \code{bookingOpensAt} (unless the student has a late-booking allowance) & \code{canBookSlot} & ``Booking has not opened yet.''\\
  Freeze window & \code{isFrozen} & ``Changes are locked Nh before the slot starts.''\\
  Capacity & \code{canBookSlot} + row lock (\code{SELECT … FOR UPDATE}) in \code{bookings.lockSlots} & ``This slot is full.''\\
  One active booking per student per assignment & \code{bookSlot} check + partial unique index \code{booking_one_active_per_student} & ``You already have a booking for this assignment --- reschedule it instead.''\\
  No overlapping demos across courses & \code{bookings.assertNoTimeClash} & ``You already have another demo booked at that time.''\\
  Change budget, reschedule & \code{canReschedule} & ``You have used all N changes.''\\
  Change budget, rebook after cancelling with none left & \code{canBookSlot} (\code{budget.lockedOut}) & ``You've used all your changes for this assignment, so you can't book again yourself. Ask your TA for a slot.''\\
  Student cancel allowed & \code{canCancelBooking} & ``Cancellations are not allowed for this assignment. Ask your TA if you can't attend.''\\
  Staff can cancel any upcoming booking, ignoring the freeze, but must give a reason & \code{staffCancelBooking} & ``Give a reason --- it's sent to the student.''\\
  Attendance only once the demo has started (resetting to pending is always allowed) & \code{canMarkAttendance} & ``You can record attendance once the demo has started.''\\
  \bottomrule
\end{xltabular}}

\textbf{Active booking} = \code{BOOKED}, \code{COMPLETED} or \code{NO_SHOW} (\code{ACTIVE_BOOKING} in \src{src/server/services/slots.ts}). \code{CANCELLED} rows are kept as history.

\textbf{Reschedule} = the old booking is cancelled (by the student) and a new one is created in one transaction.

\subsection{Change budget}\label{sec:change-budget}

One budget per student per assignment (\code{changeBudget} in \code{booking-rules.ts}, loaded by \code{bookings.loadChangeState}):

\begin{itemize}
  \item \code{allowed} = \code{maxReschedules} + the student's \code{BookingAllowance.extraChanges}.
  \item \code{used} = the number of this student's bookings for the assignment with \code{status = CANCELLED}, \code{cancelledById = student} and \code{freeCancel = false}. That covers self-cancellations and the old half of every reschedule.
  \item A reschedule needs \code{left > 0}. Cancelling is allowed while \code{allowStudentCancel} is on. With a change left it spends one, and rebooking afterwards is free. \textbf{With none left} it spends nothing (\code{cancelSpendsChange}; the booking is marked \code{freeCancel}), so \code{used} never exceeds \code{allowed}, but the student is \code{lockedOut}: they can't book again themselves until staff place them or grant an extra change. A granted change then lets them rebook and stays a real spare.
  \item Staff cancellations, slot cancellations and staff moves never count.
  \item The cancel and reschedule confirmations say exactly what they cost (\code{cancelConsequence}, \code{rescheduleConsequence}), including ``This uses your last change''.
  \item The student's page offers \ui{Reschedule} (and the slot picker, even from a \code{?reschedule=1} link) only when a move is possible.
\end{itemize}

Example with \code{maxReschedules = 2}: book A → reschedule to B (1/2) → cancel (2/2) → book C \checkmark\ → cancel (none left: still 2/2, locked out) → book $\times$ → staff grant 1 extra (2/3) → book D \checkmark\ → reschedule \checkmark\ (3/3).

\subsection[Allowances (\texttt{BookingAllowance})]{Allowances (\code{BookingAllowance})}

Staff can grant one student an exception on one assignment:
\begin{itemize}
  \item \code{extraChanges}: more changes on top of the policy.
  \item \code{lateBooking}: book, reschedule or cancel after booking has closed or before it opens.
\end{itemize}

The freeze window, capacity and draft slots still apply.

\section{Staff tools \& student requests}

{\tablefont
\begin{xltabular}{\linewidth}{@{}P{30mm} P{30mm} L@{}}
  \toprule
  \thead{Operation} & \thead{Service} & \thead{Rule}\\
  \midrule
  \endhead
  Place or move a student & \code{bookings.staffPlaceStudent} & Ignores the freeze window and booking dates. Over capacity only when ``allow over capacity'' is ticked. Refused for completed demos, past or cancelled slots, and time clashes. The old booking is cancelled by staff, so no change is used. An open booking request stays open (now pointing at the new booking) so staff can reply in their own words; the request card always shows the student's current booking. The old booking is marked \code{replacedBy = MOVED_BY_STAFF}, so the student's history says ``Moved by staff''; a student's own reschedule is marked \code{RESCHEDULED}. The student's notice gives the new and the previous time.\\
  Let a no-show book again & \code{bookings.allowRebookAfterNoShow} & Only \code{NO_SHOW} bookings whose evaluation isn't submitted. No change is used.\\
  Exception (\code{BookingAllowance}) & \code{bookings.setAllowance} & 0--20 extra changes and/or late booking. Setting both to zero removes it. The student is notified.\\
  Waitlist & \code{waitlist.joinWaitlist}, \code{notifyWaitlist} & Only while the assignment is open and the student has no booking. When a seat frees up (cancel, reschedule, staff cancel or move, more capacity, new published slots), every waiting student is notified, at most once per 30 minutes. First to book gets it. The entry is removed when they book.\\
  Student request & \code{requests.createRequest}, \code{resolveRequest} & One open request per kind (\code{BOOKING_CHANGE}, \code{MARK_QUERY}) per assignment. Mark queries only after marks are released. Goes to the student's host and the instructors (all staff if neither). Declining needs a reply.\\
  Change host & \code{slots.reassignHost} & The new host must be course staff and have no overlapping slot. Booked students are told.\\
  Slot capacity & \code{slots.updateSlotCapacity} & 1--100, never below the students already booked. More capacity alerts the waitlist.\\
  Bulk finalize & \code{evaluations.finalizeMany} & Instructor only. Same as finalizing each one.\\
  \bottomrule
\end{xltabular}}

\section[Automations]{Automations (worker, \src{src/server/services/automations.ts})}

Each job records what it did, so running it twice sends nothing twice.

{\tablefont
\begin{xltabular}{\linewidth}{@{}P{30mm} P{30mm} L@{}}
  \toprule
  \thead{Job} & \thead{When} & \thead{What}\\
  \midrule
  \endhead
  \code{announceOpenedBookings} & \code{bookingOpensAt} has passed & ``Book your demo'' to students without a booking\\
  \code{nudgeUnbookedStudents} & 48h before \code{windowEnd} (once) & ``Last chance to book'' to unbooked students\\
  \code{closeEndedAssignments} & \code{windowEnd} has passed & Assignment → \code{CLOSED}, waitlist cleared, audited with no actor\\
  \code{sendDailyAgendas} & 06:00--11:00 course time, once per local day & Email-only list of today's demos for each host (unless they opted out)\\
  \code{queueDueReminders} & 24h / 1h before & Reminder to the student (unless they opted out of reminder emails)\\
  \bottomrule
\end{xltabular}}

\section{Calendar}

\begin{itemize}
  \item \ui{Add to calendar} on a booking: Google Calendar and Outlook links that open with the demo filled in, or \route{/bookings/<id>/calendar} as an \code{.ics} file (the student's own, signed in). Each adds a copy; when the booking changes the student is emailed and can add the new time.
  \item Automatic sync is \ui{coming soon} (see the roadmap). \route{/calendar/<token>} (a personal feed) still answers for links created earlier, but the account page no longer offers it.
\end{itemize}

\section{PWA}

{\tablefont
\begin{xltabular}{\linewidth}{@{}P{30mm} P{30mm} L@{}}
  \toprule
  \thead{Piece} & \thead{Where} & \thead{Behaviour}\\
  \midrule
  \endhead
  Install & \src{src/app/manifest.ts}, \src{src/app/icons/[kind]/route.tsx}, \code{InstallHint} (dashboard) & Standalone app starting at \route{/dashboard}. Icons are generated (192, 512, maskable, badge). The dashboard offers ``Install'', or Add to Home Screen instructions on iOS. The hint can be dismissed and stays dismissed.\\
  Service worker & \src{public/sw.js}, registered by \code{ServiceWorkerRegistrar} in production (or with \code{NEXT_PUBLIC_ENABLE_SW=1}); under \code{next dev} it unregisters any worker and \code{slotty-*} caches left on the origin & Build assets are cache-first, and only responses marked \code{immutable} are stored (dev assets change at stable URLs). Dashboard, bookings, notifications, course and assignment pages are network-first with the last good copy used offline. Anything else offline shows \route{/offline}. Only GETs are cached; Server Actions and RSC fetches never are.\\
  Offline mode & \code{OfflineBanner}, \code{SubmitButton} & A banner says the page is the copy saved at a given time. Every submit button is disabled while offline. \textbf{Changes are never queued}: replaying a booking later could hit a slot that has since changed.\\
  Push & \code{PushSubscription} (tied to the login \code{Session}), \code{PushMessage} outbox, \src{src/server/push.ts} (worker) & Every in-app notification is also queued as a push for the user's devices, in the same transaction. The worker delivers pushes, drops ones older than 6h, and deletes subscriptions the push service reports gone (404/410). Opt-in is on the account page, or a prompt after booking. The permission request is never shown on page load.\\
  Privacy & \route{/logout} route, \code{ClearSavedPages} (login page) & Signing out sends \code{Clear-Site-Data: "cache", "storage"} and deletes the session, which removes that device's push subscription. The sign-in page also clears saved pages.\\
  \bottomrule
\end{xltabular}}

Limits: push arrives even when the app is closed, but the device needs a connection. iOS supports push only after Add to Home Screen (16.4+). Push and service workers need HTTPS (localhost is exempt).

\section[Scheduling rules]{Scheduling rules (\src{src/domain/slots.ts}, \src{src/server/services/slots.ts})}

\begin{itemize}
  \item An availability block is split into \code{slotDurationMin} slots with \code{bufferMin} gaps and clipped to the demo window. A leftover shorter than one slot is dropped.
  \item Slots can be added for several days at once (same hours each day).
  \item A host can't have overlapping non-cancelled slots in \textbf{any} assignment. Times that clash are \textbf{skipped} and reported; if every time clashes, nothing is added.
  \item The host must be a TA or instructor of the course, and the venue must belong to the course. \textbf{TAs add slots only for themselves}; instructors (and admins) can add them for any staff member, and only they can hand slots to another host.
  \item Slots added before booking opens are hidden (\code{DRAFT}). Slots added while booking is open can be booked straight away.
  \item Only slots with no bookings at all can be deleted. Cancelling a slot releases its bookings and notifies the students.
  \item A venue can't be deleted while active slots use it.
\end{itemize}

\section[Assignment rules]{Assignment rules (\src{src/server/services/assignments.ts})}

\begin{itemize}
  \item Rubric points must add up to \textbf{exactly} \code{maxMarks} (\code{checkRubric} in \src{src/domain/evaluation.ts}): fewer would make full marks impossible (6 + 3 of 10 tops out at 9), more would let totals pass the maximum. No rubric means marking one total. Edits check this only when the rubric or max marks change, so an older assignment whose locked rubric doesn't add up can still be edited. The rubric can't change once any score has been saved.
  \item \code{maxMarks} can't go below the highest total already given.
  \item The demo window can't be changed so that it no longer covers booked slots. Unbooked slots left outside the window are reported so staff can delete them.
  \item Slot length, break and capacity only apply to slots added later.
  \item Changing the freeze window, the number of changes allowed, or \code{allowStudentCancel} notifies students with a booking (\code{assignment.rules_changed}).
  \item Opening booking (\code{publishAssignment}, \ui{Open booking} in the UI) needs at least one non-cancelled slot, because every student is notified. Reopening a closed assignment uses the same action. If \code{bookingOpensAt} is in the future, students are told the opening time (\code{assignment.opens_soon}), and the worker sends ``Book your demo'' when it arrives.
  \item Closing stops bookings and changes. Existing bookings and marking carry on.
  \item An assignment with any booking can't be deleted. Close it instead.
\end{itemize}

\section[Evaluation states]{Evaluation states (\src{src/domain/evaluation.ts})}

\begin{figure}[htbp]
  \centering
  % Evaluation state machine — mirrors `nextStatus` / `canTransition` in src/domain/evaluation.ts.
\begin{tikzpicture}[node distance=10mm and 34mm]
  \node[state] (draft) {DRAFT};
  \node[state, right=of draft] (sub) {SUBMITTED};
  \node[terminal, right=of sub] (fin) {FINALIZED};
  \node[state, below=16mm of sub] (ret) {RETURNED};
  \draw[trans] ([yshift=1.5mm]draft.east) -- node[translabel, above, align=center]{submit\\\itshape review, or held 24\,h} ([yshift=1.5mm]sub.west);
  \draw[trans] ([yshift=-1.5mm]sub.west) -- node[translabel, below]{withdraw (held)} ([yshift=-1.5mm]draft.east);
  \draw[trans] (sub) -- node[translabel, above]{finalize} node[translabel, below, align=center]{\itshape or release\\\itshape at releaseAt} (fin);
  \draw[trans] (sub.south west) ++(4mm,0) -- node[translabel, left, align=right]{return\\\itshape comment required} ++(0,-16mm+7mm);
  \draw[trans] (fin.south) |- node[translabel, pos=.25, right, align=left]{unlock\\\itshape reason required} (ret.east);
  \draw[trans] (ret.north east) ++(-4mm,0) -- node[translabel, right]{submit} ++(0,16mm-7mm);
  \draw[trans] (draft.north) to[out=50, in=130] node[translabel, above, align=center]{submit\\\itshape by an instructor} (fin.north);
  \draw[trans] (draft.south west) ++(2mm,0) arc[start angle=0, end angle=-300, radius=2.6mm] node[translabel, left=1mm, pos=.5]{save};
  \draw[trans] (ret.south west) ++(2mm,0) arc[start angle=0, end angle=-300, radius=2.6mm] node[translabel, left=1mm, pos=.5]{save};
\end{tikzpicture}

  \caption{Evaluation states (\src{src/domain/evaluation.ts}).}
  \label{fig:evaluation-states}
\end{figure}


\begin{itemize}
  \item \code{submit} with no instructor in the course → \code{FINALIZED} directly.
  \item Only \code{DRAFT} and \code{RETURNED} can be edited. Every rubric row must be scored to submit, and the total must be between 0 and \code{maxMarks}.
  \item What an \textbf{instructor} (or admin) submits is final: they are the reviewer, so it skips \code{SUBMITTED}.
  \item Marking \textbf{before the demo starts} needs a reason (\code{earlyMarkReason}, audited as \code{evaluation.early}).
  \item \textbf{Recorded marks pin the booking.} A booking with marks (any total, or submitted/final) can't be cancelled or moved by the student or staff, and its slot can't be cancelled. The student's page hides \ui{Cancel booking} and \ui{Reschedule} and says why (\code{isMyBookingMarked}). Staff \ui{Clear marks} first (draft only, with a reason, audited as \code{evaluation.clear}; submitted or final marks are unlocked first).
  \item The marking sheet autosaves drafts (\code{saveDraft}); the evaluation is created on the first save, never just by viewing. A score for a booked student whose demo has started records them as present. A no-show can't be marked until the no-show is undone.
  \item Marking a student \textbf{with no booking} needs a reason. It's stored on the evaluation (\code{noBookingReason}), flagged in lists, and audited as \code{evaluation.no_booking}.
  \item A total override needs a note. Attendance locks once the evaluation is \code{SUBMITTED} or \code{FINALIZED}.
  \item Only instructors review. Unlocking needs a reason (in a course without an instructor, the TA can unlock).
  \item Students see released results only, never drafts or \code{privateNotes}. Unlocking keeps a copy of what the student was shown (\code{Evaluation.released}); until the marks are released again the student keeps seeing it, marked as being corrected, together with any marks question. Releasing clears the copy.
  \item Releasing marks (when submitting is final) asks for confirmation first, and every message uses one word for it: ``released''. Otherwise marks are ``sent for review''.
  \item The marking sheet shows the total as ``--'' while any mark is invalid (non-numeric or outside 0 to the row's maximum).
\end{itemize}

\section[Access]{Access (\src{src/server/services/access.ts})}

\begin{itemize}
  \item Every service takes an \code{Actor} and checks the role itself: \code{assertCourseRole}, \code{assertStaff}, \code{assertAdmin}, \code{assertPlatform}.
  \item \textbf{One kind of staff per person; students are never instructors.} A TA can be a student elsewhere, but an instructor is never a TA or a student in any course, and a TA or student is never an instructor (\code{canTakeRole}, \src{src/domain/roles.ts}). Imports skip such rows with a reason; other paths refuse.
  \item \textbf{Nobody switches between student and staff in a course} (\code{canTakeRole}, ADR-025). A student there can't be made a TA or instructor, and staff can't be made students, by anyone, admins included. Removing and re-adding doesn't get round it: someone who booked or was marked in the course can't join it as staff, and someone who hosted or marked demos in it can't join as a student (\code{pastRolesHere}).
  \item \textbf{Only admins name instructors} (\code{canManageInstructors}), from the console (\ui{New course}, \ui{Assign staff}, \ui{Remove staff member}). Adding, changing, demoting or removing an instructor from inside a course is refused for everyone.
  \item \textbf{Only instructors manage TAs} (\code{canManageTAs}, ADR-025). TAs add and remove students only. In a course without an instructor, co-TAs are added and removed by an admin from the console. \ui{Add someone} offers \ui{TA} only to instructors, and the import help and template follow. Re-importing a TA unchanged isn't a change and is allowed.
  \item \textbf{Self-service courses are run by a TA} (\code{canCreateCourse}). \route{/courses/new} always makes the creator the course's TA; there is no role choice, and the server ignores one. It is open to TAs and to new staff who aren't in any course yet; students (in a course, TA nowhere) and instructors (who can't also be TAs) are refused, and \ui{New course} is hidden for them.
  \item \textbf{The instructor manages the course; a TA only when there is none} (\code{canManageCourse}, \code{assertCourseManager}). Course details and archiving need it; \ui{Settings} is hidden otherwise.
  \item \textbf{Admins read courses but never change them} (\code{canChangeCourseData}, ADR-025). They open any course's pages as if they were its instructor (``Viewing as admin'', with a read-only note), but every change is refused: marks, bookings, slots, requests, venues, assignments, roster, settings and archiving. \code{assertCourseWritable(actor, course)} applies it, with the archived check, to every mutation. They are never enrolled either: roster imports skip admin emails, \ui{Assign staff} and \ui{Add someone} refuse them, they can't create a course for themselves (they create it from the console and name its instructor or TA), and a course member can't be made admin.
  \item \textbf{Access before data.} A roster import or preview checks the caller's access before it screens the rows, since skip reasons name people's other courses. Staff-only counts and lists (\code{listCourseStaff}, \code{openRequestCount}, \code{pendingReviewCount}, \code{markingSheetFocus}) check access themselves rather than relying on the page.
  \item Nobody can remove themselves, and nobody can be demoted from staff to student (above). Adding someone already in the course with a blank section keeps their section (a blank never wipes it), and says they're already there.
  \item \textbf{Archived courses are read-only.} Mutations pass \code{{ write: true }} to \code{assertCourseRole} / \code{assertStaff} / \code{assertAssignmentRole}, or call \code{assertCourseWritable}. Only restoring the course is allowed.
  \item Removing a student cancels their upcoming bookings and notifies them. Removing staff is refused while they host upcoming slots.
  \item Admins disabling an account can also release that user's upcoming bookings.
  \item \textbf{Platform accounts are only platform accounts} (\code{User.platformAccess}, \code{canGrantPlatform}). Usage insights and feedback need \code{assertPlatform}, which refuses admins (and any record with both flags). Access is granted only by \src{src/jobs/platform-access.ts}, which refuses admins, course members and inactive accounts; \src{src/jobs/create-admin.ts} refuses platform accounts. Platform accounts are never enrolled (roster imports skip them, \ui{Assign staff} and course creation refuse them, \code{assertCanJoinCourses}), and the admin console doesn't list, count, open or change them; their activity and \code{platform.*} entries are left out of the admin audit log.
\end{itemize}

\subsection*{Permission matrix}

What each role may do in a course (\code{assertCourseRole} and the rules above). ``Own'' means
the actor's own bookings, slots or requests. Admins and platform accounts are never enrolled;
platform accounts see no course at all.

{\tablefont
\begin{xltabular}{\linewidth}{@{}L c c c c@{}}
  \toprule
  \thead{Action} & \thead{Instructor} & \thead{TA} & \thead{Student} & \thead{Admin}\\
  \midrule
  Open the course's pages, roster, slots, marks and requests & yes & yes & own view & read\\
  Book, reschedule, cancel; join the waitlist; ask a question & -- & -- & own & --\\
  Course details and archiving & yes & without instructor & -- & console\\
  Add or remove students & yes & yes & -- & --\\
  Add, promote, demote or remove TAs & yes & -- & -- & console\\
  Add, change or remove instructors & -- & -- & -- & console\\
  Switch someone between student and staff & -- & -- & -- & --\\
  Assignments, venues, and slots for any host & yes & yes, slots for self & -- & --\\
  Mark attendance, place or cancel students' bookings, grant exceptions & yes & yes & -- & --\\
  Mark and submit evaluations & yes & yes & -- & --\\
  Review, finalize and unlock marks & yes & unlock without instructor & -- & --\\
  Answer students' requests & yes & yes & -- & --\\
  \bottomrule
\end{xltabular}}

\section{Accounts}

The account page (\route{/account}) shows only what applies to the user's roles: name, clock (\code{timeFormat}: 12 or 24 hour), password (checks the current one and signs out other devices), ``sign out other devices''; reminder emails (\code{emailReminders}) for students, the morning agenda (\code{emailAgenda}) for staff, and device notifications for both. Admins see none of the demo settings. Confirmation, change and marks emails can't be turned off.

{\tablefont
\begin{xltabular}{\linewidth}{@{}P{30mm} L@{}}
  \toprule
  \thead{Item} & \thead{Value}\\
  \midrule
  \endhead
  Invite link & 14 days. Re-inviting invalidates older links.\\
  Password reset link & 2 hours. Resetting ends every session.\\
  Session & 30 days, stored in the DB. Disabled users are signed out.\\
  Unknown email on reset & Silent success, so nobody can probe which emails exist.\\
  \bottomrule
\end{xltabular}}

\section[Notifications]{Notifications (\src{src/server/services/notify.ts})}

Each notification creates an in-app entry and, for people who have activated their account, an email (outbox). Invited people get only their invite email; everything else waits in-app for when they join. Both are written in the \textbf{same transaction} as the change. The worker (\code{bun run worker}) sends emails and queues reminders.

How much people are told (\src{src/domain/notifications.ts}, ADR-018):
\begin{itemize}
  \item \textbf{Not about their own actions.} A student's own booking, reschedule or cancel creates no in-app item or push (the page already says so). One confirmation email for where they end up goes out \code{SELF_EMAIL_DELAY_MINUTES} (10) later.
  \item \textbf{Supersede, don't stack.} Items about the same thing share a \code{groupKey} (one student's booking for an assignment; one evaluation's marks). A new one removes unread items and unsent emails with the same key from the last \code{SUPERSEDE_MINUTES} (15). Booking emails list what they replaced (``Earlier: …'').
  \item \textbf{Coalesced marks.} Marks emails wait 5 minutes. Corrected marks released before the student saw the first release (\code{unseenEarlier}) become one ``Marks released'' item with the final marks; later corrections are ``Marks updated'' with the old total.
  \item \textbf{The bell counts what needs the person} (\code{needsAction}): \code{course.enrolled}, \code{slots.added}, \code{assignment.rules_changed}, \code{slot.venue_changed} and \code{agenda.daily} are listed but not counted. The list is grouped by assignment (\code{groupByAssignment}).
  \item Times in messages name the zone as the UI does (\code{tzLabel}: ``Pakistan time'').
\end{itemize}

{\tablefont
\begin{xltabular}{\linewidth}{@{}P{30mm} P{30mm} L@{}}
  \toprule
  \thead{Type} & \thead{Sent to} & \thead{When}\\
  \midrule
  \endhead
  \code{course.enrolled} & existing user & Added to a course (new users get an invite email instead). Students are also told which demos are already open.\\
  \code{course.removed} & student & Removed from a course while holding upcoming bookings\\
  \code{assignment.opens_soon} & all students & Published with a future \code{bookingOpensAt}\\
  \code{assignment.opening_moved} & all students & The future opening time was changed\\
  \code{slots.added} & students without a booking (not waitlisted) & Slots added while booking is open\\
  \code{assignment.published} & students without a booking & Booking is open: on publish or reopen, or by the worker at \code{bookingOpensAt}\\
  \code{assignment.rules_changed} & booked students & Freeze window, changes allowed or cancel permission changed\\
  \code{booking.confirmed} & student (email only, delayed) & Booked\\
  \code{booking.rescheduled} & student (email only, delayed) & Rescheduled\\
  \code{booking.cancelled} & student (email only, delayed) & Student cancelled\\
  \code{booking.cancelled_by_staff} & student & Staff cancelled their booking\\
  \code{slot.cancelled} & booked students & Staff cancelled the slot\\
  \code{slot.venue_changed} & booked students & Venue changed\\
  \code{booking.reminder} & student & 24h and 1h before (worker)\\
  \code{booking.no_show} & student & Marked as no-show\\
  \code{booking.placed_by_staff} & student & Staff booked or moved them, with the previous time\\
  \code{booking.rebook_allowed} & student & Staff cleared a no-show\\
  \code{booking.allowance} & student & Exception granted\\
  \code{waitlist.slot_available} & waiting students & A seat freed up\\
  \code{slot.host_changed} & booked students & Slot handed to another host\\
  \code{request.created} & host + instructors & Student sent a request\\
  \code{request.answered} & student & Staff handled or declined it\\
  \code{assignment.book_soon} & unbooked students & 48h before the demo window ends\\
  \code{agenda.daily} & host (email only) & Morning list of today's demos\\
  \code{evaluation.returned} & evaluator (TA) & Instructor returned an evaluation\\
  \code{evaluation.finalized} & student & Marks released; the body has the total, each rubric row and the feedback\\
  \code{evaluation.updated} & student & Corrected marks released after an unlock, with the old total (only once the first release was seen)\\
  \bottomrule
\end{xltabular}}

\section[Usage data and feedback]{Usage data and feedback (\src{src/domain/analytics.ts})}\label{sec:rules-usage}

{\tablefont
\begin{xltabular}{\linewidth}{@{}P{38mm} L@{}}
  \toprule
  \thead{Rule} & \thead{Behaviour}\\
  \midrule
  \endhead
  Collection switch & Usage events are recorded only where \envvar{ANALYTICS_SECRET} is set (production), and never for platform accounts. Development, integration and end-to-end runs record nothing. Feedback and surveys work everywhere.\\
  Consent & \code{User.shareUsageData}, on by default (opt-out, explained at \route{/privacy}). Off: only \code{client_error} events are kept, and without the visitor key (\code{allowedWithoutConsent}).\\
  Event catalog & \code{clientEventSchema}: \code{page_view}; \code{web_vital} (LCP, INP, CLS, FCP, TTFB, value 0--600\,000); \code{funnel} (flow \code{booking}, \code{reschedule} or \code{marking}; step \code{view}, \code{submit}, \code{done} or \code{error}); \code{feature} (\code{add_to_calendar}, \code{push_enabled}, \code{app_installed}); \code{client_error}. Anything else, and any event with an unknown field value, is dropped one by one. At most 50 events per request and 240 per person per minute.\\
  No identifiers & Paths are stored as route templates (\code{routeTemplate}: any segment that isn't a plain word becomes \code{:id}; query strings dropped). The person is a keyed hash of their user id (\code{pseudonym}). The device is reduced to mobile, tablet or desktop (\code{deviceClass}); neither IP address nor user agent is stored. Error messages lose emails, links, long numbers and token-like strings (\code{scrubErrorMessage}), and are cut to 200 characters.\\
  Times & A client timestamp is used only if it is within the last hour (\code{eventTime}), else the server's time.\\
  Retention & Raw events 90 days, daily totals 400 days (about 13 months), deleted hourly by the worker (\code{retentionCutoffs}, \code{rollupAndPrune}).\\
  Feedback & A message (up to 2000 characters) is required, the rating 1--5 is optional; at most 5 per person per hour (\code{canSendFeedback}, counted in memory so nothing links a message to its sender). Stored without the user, with the time rounded to the hour.\\
  Surveys & One question each after a booking, a reschedule and a mark submission (\code{SURVEY_PROMPTS}); never for admins or platform accounts. Each question at most once every 90 days and no survey at all within 14 days of another, answered or not (\code{canPromptSurvey}), so one demo never brings two. An answer (1--5, or dismissed) counts only within 24 hours of being asked and only once (\code{canAnswerSurvey}). Responses are stored without the user, dated to the day; \code{SurveyPrompt} keeps only when each person was asked and whether they answered.\\
  \bottomrule
\end{xltabular}}

\section{Time}

Times are stored in UTC. They're shown and entered in the \textbf{course timezone} (\src{src/lib/time.ts}: \code{fmt}, \code{fmtTimeRange}, \code{fromLocalInput}). Courses can use Pakistan time only for now (\code{SUPPORTED_TIMEZONES}). Each user picks a 12- or 24-hour clock; \code{fmt} reads it for the current request (\src{src/server/clock.ts}), while data formats (CSV, form values) stay 24-hour via \code{fmtData}. Emails sent by the worker use the 12-hour clock.

\section{Audit}

Moderation-relevant changes go through \code{audit(tx, actor, …)}: courses, roster, assignments, availability, slot cancellation and venue changes, attendance, staff cancellations, evaluation transitions, and admin actions. Platform actions (access grants, feedback status) are recorded as \code{platform.*}, outside the admin audit log.
